% Chapter 8: How Do You Know
% Scene function: Meri asks a craft question that becomes personal.
% Tert can't distinguish his work from Gordon's grief. Catches
% himself using Management's language for the second time.

\chapter{Trust the Process}

Meri is at my desk when I get back from a shift.

This is not standard. Nobody sits at someone else's desk. There was an incident in 2019 involving a borrowed pen that generated a seven-page Compliance memo on workspace sovereignty. The memo resolved the situation by existing. Meri has either not read it or has decided that her question is worth the breach.

``How do you know when it is working?'' she says.

I set my bag down. I do not ask her to move. Escalating a minor boundary issue into a conversation about boundaries would take longer than whatever Meri actually wants to talk about.

``When what is working?''

``The ambient conditioning. Phase 1. You shift the thermostat, you adjust the timing of the house sounds, you introduce sub-threshold variations. How do you know any of it is landing?''

Meri has been reading. Not the Operations Manual. Something useful, which narrows the options considerably. Shift reports, maybe, or the Phase 1 technical supplement that Compliance published in 2011 and nobody acknowledged.

``Behavior changes,'' I say. ``Sleep is the first indicator. Then attention span. Then routine drift: small changes in when the subject eats, when they go to bed, whether they finish tasks they start. You are looking for erosion, not collapse. If the subject collapses, you escalated too fast, and you will be filling out paperwork for the rest of the quarter.''

``What does erosion look like?''

``Sunday night my subject made a cup of tea and did not drink it.''

Meri considers this. ``That could be anything. That could be grief.''

``Yes.''

``So how do you know it was you?''

This is the question. Not the pipeline question from the break room, which was abstract and philosophical. This one is concrete. How do I know my work is having an effect, as opposed to the subject's grief doing it for me?

The answer is that I do not know. I have no way of distinguishing between ambient unease caused by my interventions and ambient unease caused by being a fifty-five-year-old widower alone in the house where his wife died. The symptoms are identical. The subject's sleeplessness could be Phase 1 taking hold or it could be Tuesday.

``You learn to trust the process,'' I say, which is the second time I have used one of Management's phrases to avoid answering a question from Meri, and I like it less than I did the first time.

Meri looks at me for a moment. She does not write this one down.

``Thank you,'' she says, and leaves.
